\section{Introduction}
\label{sec:intro}

Classical Arabic poetry is written in lines called \emph{bayts}. Each bayt has two halves, or hemistichs, and every
hemistich must follow one of sixteen quantitative meters, the \emph{bu\d{h}\=ur}, codified in the eighth century by
al-Khal\={\i}l ibn A\d{h}mad~\cite{frolov2000classical}. Every bayt of a poem must also end on the same rhyme letter,
the \emph{raw\={\i}}. These rules are strict. A reader trained in the tradition hears a missing syllable at once, and
a teacher of Arabic literature will reject a poem with one broken hemistich. A system that claims to write classical
Arabic verse therefore has a clear and checkable target.

Large language models (LLMs) have made such systems easy to build. Shaer~\cite{abbas2026shaer} fine-tunes the Arabic
model Yehia-7B~\cite{yehia2025} with QLoRA~\cite{dettmers2023qlora} on 116,032 poems from the Ashaar
corpus~\cite{alyafeai2023ashaar}. A prompt names the meter, its sub-form, the number of hemistichs and a topic
description, and the model writes the poem. The authors report a base-meter accuracy of 95.17\% and an exact length
accuracy of 83.40\% on 3,481 held-out prompts, far above general Arabic models such as Fanar~\cite{fanar2026fanar2}.

The meter numbers come from a bidirectional LSTM (BiLSTM) classifier~\cite{abandah2022classifying} that reads one
bayt at a time and outputs the most likely meter. A classifier of this kind is trained to tell the sixteen meters
apart. It is not trained to say ``this is not a meter at all''. A bayt that is one syllable short of al-K\=amil is
still closer to al-K\=amil than to any other meter, so the classifier will label it al-K\=amil with high confidence.
High accuracy under this metric may therefore hide many unmetrical lines. Rhyme is not checked at all.

This paper is an empirical study of that gap. We first \emph{reproduce} Shaer end to end on Google Colab GPUs, and
then \emph{extend} its evaluation with a formal verifier of Arabic prosody, which we use to measure how often the
poems are actually correct and to guide generation and training. We ask three research questions:
\begin{itemize}
\item \textbf{RQ1 (reproduction).} Can Shaer's published results be reproduced from its released code, data and
adapter, and where do the numbers differ?
\item \textbf{RQ2 (measurement).} How often are the generated poems strictly valid in meter, rhyme and length, and how
far does the classifier-based metric overstate this?
\item \textbf{RQ3 (improvement).} Can the verifier raise strict validity when it is used at decoding time or to select
training data, without hurting the quality of the poems?
\end{itemize}

\textbf{Contributions.} The paper makes the following contributions.
\begin{enumerate}
\item \emph{An independent reproduction} of Shaer (released poems, released adapter, untuned base model and a full
re-training), reported next to the published numbers with an explanation of every difference
(Section~\ref{sec:res-repro}).
\item \emph{A formal, open verifier of classical Arabic prosody} that derives 1,867 hemistich patterns of 16 meters
from the classical rules, scans undiacritised or partly diacritised text with a pronunciation prior, and checks
rhyme (raw\={\i}, its vowel, \emph{ridf}, \emph{ta'\={\i}s}) and length (Section~\ref{sec:m1}).
\item \emph{A strict re-scoring} of the reproduced systems and of the human poems, showing that strict validity is far
below what the classifier metric suggests and that it falls with poem length (Section~\ref{sec:results}).
\item \emph{A robustness analysis} with nested logistic models, nested definitions of validity, alternative settings
and control samples (Section~\ref{sec:robust}).
\item \emph{Three pre-registered extensions}: verifier-guided hemistich search compared with best-of-$N$ sampling at
equal compute, self-training on verifier-approved poems, and a quality study with an LLM judge and blinded human
raters. Their designs and predictions were fixed before the runs (Sections~\ref{sec:m4}--\ref{sec:m6}).
\end{enumerate}

The rest of the paper is organised as follows. Section~\ref{sec:lit} reviews related work and states the hypotheses.
Section~\ref{sec:background} introduces the prosody needed to read the paper. Section~\ref{sec:methods} describes the
techniques M0--M6 with examples from the data. Sections~\ref{sec:design} and~\ref{sec:setup} give the research
design and the analysis setup. Sections~\ref{sec:results} and~\ref{sec:robust} report the results and the robustness
checks. Section~\ref{sec:case} discusses a practical case, Section~\ref{sec:ethics} ethics and limitations, and
Section~\ref{sec:conclusion} concludes.
